\documentclass[11pt]{article}
\usepackage[a4paper,margin=1in]{geometry}
\usepackage[T1]{fontenc}
\usepackage[utf8]{inputenc}
\usepackage{lmodern}
\usepackage{amsmath,amssymb,amsthm,mathtools,bm}
\usepackage{mathrsfs}
\usepackage{hyperref}
\usepackage{enumitem}
\usepackage{microtype}
\usepackage{xcolor}
\usepackage{booktabs}
\usepackage{dsfont}
\hypersetup{colorlinks=true,linkcolor=blue,citecolor=blue,urlcolor=blue}
\setlength{\parskip}{0.4em}
\setlength{\parindent}{0pt}

\newtheorem{proposition}{Proposition}
\newtheorem{remark}{Remark}
\newtheorem{definition}{Definition}

\newcommand{\Tr}{\operatorname{Tr}}
\newcommand{\e}{\mathrm{e}}
\newcommand{\ii}{\mathrm{i}}
\newcommand{\dd}{\mathrm{d}}
\newcommand{\1}{\mathds{1}}
\newcommand{\HA}{\mathcal{H}_A}
\newcommand{\HB}{\mathcal{H}_B}
\newcommand{\HH}{\mathcal{H}}
\newcommand{\cD}{\mathcal{D}}
\newcommand{\up}{\uparrow}
\newcommand{\dn}{\downarrow}
\newcommand{\bh}{\bm{h}}
\newcommand{\bI}{\bm{I}}
\newcommand{\br}[1]{\left(#1\right)}
\newcommand{\sq}[1]{\left[#1\right]}
\newcommand{\cb}[1]{\left\{#1\right\}}

\title{From the Fermionic Graded-Swap Representation of $\Tr_A(\rho_A^n)$\\
to the Free Replica Green's Function, Perturbation Theory,\\
and the Shifted Expansion}
\author{}
\date{}

\begin{document}
\maketitle

\begin{abstract}
We present a self-contained derivation of the fermionic replica construction
for the $n$-th Renyi entropy of a spatial subregion $A$, starting from
\[
Z_n[A]\equiv \Tr_A(\rho_A^n),
\qquad
\rho=\frac{\e^{-\beta H}}{Z_1},
\qquad
Z_1=\Tr_{\HH}\e^{-\beta H}.
\]
The key correction relative to the naive bosonic-looking construction is that
the operator inserted on subsystem $A$ is the \emph{graded} cyclic
permutation, not the plain cyclic permutation. This graded structure changes
the temporal gluing on $A$, shifts the replica-sector phases by a half-step,
and therefore changes the free replica Green's function, the free replica
partition function, and the perturbative coefficients.

We derive the coherent-state path-integral representation with the corrected
mixed temporal boundary conditions, introduce the graded replica shift matrix,
and diagonalize it by a replica Fourier transform with momenta
\[
q_m=\frac{(2m+1)\pi}{n}, \qquad m=0,\dots,n-1.
\]
In a convention without an extra minus sign in front of the time-ordering
symbol,
\[
\bar G \equiv \langle T_\tau c c^\dagger\rangle_0,
\]
the equation of motion carries a positive delta function. We derive the exact
piecewise free Green's function in one graded replica sector, discuss the loss
of imaginary-time translation invariance, and analyze the singularity
criterion. We then derive the determinant perturbation expansion and the
shifted expansion used in connected-determinant / diagrammatic Monte Carlo
implementations.

We also separate clearly the twisted unnormalized object
\[
\widetilde Z_n[A]
\equiv
\Tr_{\HH^{\otimes n}}\sq{S_A^{\mathrm{gr}}\br{\e^{-\beta H}}^{\otimes n}}
\]
from the normalized quantity
\[
Z_n[A]=\Tr_A(\rho_A^n)=\frac{\widetilde Z_n[A]}{Z_1^n}.
\]
This distinction is essential at first order: the twisted replica theory
computes the derivative of $\ln \widetilde Z_n[A]$, whereas the normalized
object relevant for the Renyi entropy requires subtracting $n \ln Z_1$.
\end{abstract}

\tableofcontents

\section{Setup and normalization convention}

Let the total Hilbert space factorize as
\begin{equation}
\HH=\HA\otimes \HB,
\label{eq:Hfactor}
\end{equation}
where $A$ is the subsystem of interest and $B=\bar A$ is its complement. The
thermal density matrix is
\begin{equation}
\rho=\frac{\e^{-\beta H}}{Z_1},
\qquad
Z_1=\Tr_{\HH}\e^{-\beta H},
\qquad
\rho_A=\Tr_B\rho.
\label{eq:rho-def}
\end{equation}
We define the normalized replica partition function by
\begin{equation}
Z_n[A]\equiv \Tr_A(\rho_A^n).
\label{eq:Zn-normalized}
\end{equation}
Hence the Renyi entropy is
\begin{equation}
S_A^{(n)}=\frac{1}{1-n}\ln Z_n[A].
\label{eq:renyi-normalized}
\end{equation}

It is often useful to introduce the twisted but unnormalized quantity
\begin{equation}
\widetilde Z_n[A]\equiv
\Tr_{\HH^{\otimes n}}\sq{S_A^{\mathrm{gr}}\br{\e^{-\beta H}}^{\otimes n}},
\label{eq:Zn-unnormalized}
\end{equation}
where $S_A^{\mathrm{gr}}$ is the fermionic graded cyclic permutation operator
to be defined below. The relation between the two objects is
\begin{equation}
Z_n[A]=\frac{\widetilde Z_n[A]}{Z_1^n},
\qquad
\ln Z_n[A]=\ln \widetilde Z_n[A]-n\ln Z_1.
\label{eq:Zn-tilde-relation}
\end{equation}

\begin{remark}
Equations \eqref{eq:renyi-normalized} and \eqref{eq:Zn-tilde-relation} are
equivalent ways of organizing the same physics. The important point is that
one must use one convention consistently. In particular, if one computes
\[
\widetilde Z_n[A]
=
\Tr_{\HH^{\otimes n}}\sq{S_A^{\mathrm{gr}}\br{\e^{-\beta H}}^{\otimes n}},
\]
then the normalized quantity entering the entropy is obtained only after
subtracting $n \ln Z_1$.
\end{remark}

\section{Replica Hilbert space and the graded cyclic permutation operator}

On $n$ replicas, the Hilbert space is
\begin{equation}
\HH^{\otimes n}=(\HA\otimes \HB)^{\otimes n}
\cong
\HA^{\otimes n}\otimes \HB^{\otimes n}.
\label{eq:replica-space}
\end{equation}
Choose homogeneous basis vectors
$\{\lvert a\rangle\}\subset\HA$ and $\{\lvert b\rangle\}\subset\HB$, where
``homogeneous'' means that each $\lvert a\rangle$ has a definite fermion
parity $p(a)\in\{0,1\}$. A replicated basis vector is
\begin{equation}
\lvert a_1,b_1;\dots;a_n,b_n\rangle
\equiv
\bigotimes_{\alpha=1}^n
\br{\lvert a_\alpha\rangle_A\otimes \lvert b_\alpha\rangle_B}.
\label{eq:replicated-basis}
\end{equation}

\begin{definition}[Graded cyclic permutation on subsystem $A$]
The fermionic replica construction uses the \emph{graded} cyclic permutation
operator $S_A^{\mathrm{gr}}$, not the plain cyclic permutation. On homogeneous
basis vectors it acts as
\begin{equation}
S_A^{\mathrm{gr}}
\lvert a_1,b_1;\dots;a_n,b_n\rangle
=
(-1)^{p(a_1)\sum_{\beta=2}^n p(a_\beta)}
\lvert a_2,b_1;a_3,b_2;\dots;a_1,b_n\rangle.
\label{eq:Sagr-action}
\end{equation}
Equivalently,
\begin{equation}
S_A^{\mathrm{gr}}
=
\sum_{a_1,\dots,a_n}^{\mathrm{hom.}}
(-1)^{p(a_1)\sum_{\beta=2}^n p(a_\beta)}
\lvert a_2,\dots,a_n,a_1\rangle_A
\,{}_A\langle a_1,\dots,a_n|
\otimes \1_B^{\otimes n}.
\label{eq:Sagr-explicit}
\end{equation}
\end{definition}

\begin{remark}
Equation \eqref{eq:Sagr-action} is the standard Koszul sign associated with a
cyclic permutation in a fermionic tensor product. This graded sign is the
entire reason that the replica-sector phases are shifted from
$2\pi m/n$ to $(2m+1)\pi/n$.
\end{remark}

\section{Graded-swap representation of $\Tr_A(\rho_A^n)$}

The matrix elements of the reduced density matrix are
\begin{equation}
(\rho_A)_{aa'}
\equiv
\langle a\lvert \rho_A\rvert a'\rangle
=
\sum_b \langle a,b\lvert \rho \rvert a',b\rangle.
\label{eq:rhoA-components}
\end{equation}
Therefore
\begin{align}
Z_n[A]
&=\Tr_A(\rho_A^n)
\notag\\
&=
\sum_{a_1,\dots,a_n}
(\rho_A)_{a_1a_2}
(\rho_A)_{a_2a_3}
\cdots
(\rho_A)_{a_n a_1}
\notag\\
&=
\sum_{a_1,\dots,a_n}
\sum_{b_1,\dots,b_n}
\langle a_1,b_1\lvert \rho \rvert a_2,b_1\rangle
\langle a_2,b_2\lvert \rho \rvert a_3,b_2\rangle
\cdots
\langle a_n,b_n\lvert \rho \rvert a_1,b_n\rangle.
\label{eq:Zn-component}
\end{align}
This makes the sewing pattern manifest:
\begin{itemize}[leftmargin=2em]
\item the $B$ index closes within the same replica,
\item the $A$ index closes cyclically across replicas.
\end{itemize}

For fermions, the corresponding operator identity uses the graded swap:
\begin{proposition}[Graded-swap representation]
For any fermionic density matrix $\rho$ on $\HH=\HA\otimes \HB$,
\begin{equation}
\Tr_A(\rho_A^n)
=
\frac{1}{Z_1^n}
\Tr_{\HH^{\otimes n}}
\sq{
S_A^{\mathrm{gr}}
\br{\e^{-\beta H}}^{\otimes n}
}
=
\frac{\widetilde Z_n[A]}{Z_1^n}.
\label{eq:graded-swap-identity}
\end{equation}
\end{proposition}

\begin{remark}
If one replaced $S_A^{\mathrm{gr}}$ by the ordinary cyclic permutation, the
full-system limit $A=\HH$ would be wrong already in the free theory: for
$n=2$ one would obtain a bosonic-looking periodic sector instead of the
correct fermionic half-shifted sectors $\pm \ii$.
\end{remark}

\section{Fermionic coherent states and the graded path integral}

Let $c_i,c_i^\dagger$ denote fermionic annihilation and creation operators,
where $i$ collects all one-particle labels. For Grassmann variables
$\eta_i,\bar\eta_i$, define the coherent states
\begin{equation}
\lvert \eta\rangle
\equiv
\exp\br{-\sum_i \eta_i c_i^\dagger}\lvert 0\rangle,
\qquad
\langle \bar\eta\rvert
\equiv
\langle 0\rvert \exp\br{-\sum_i c_i \bar\eta_i}.
\label{eq:coherent-state-def}
\end{equation}
They satisfy
\begin{equation}
c_i\lvert\eta\rangle=\eta_i\lvert\eta\rangle,
\qquad
\langle\bar\eta\rvert c_i^\dagger=\langle\bar\eta\rvert \bar\eta_i,
\label{eq:coherent-eigen}
\end{equation}
with overlap
\begin{equation}
\langle \bar\eta\vert \eta'\rangle
=
\exp\br{\sum_i \bar\eta_i \eta'_i}.
\label{eq:coherent-overlap}
\end{equation}
The identity resolution is
\begin{equation}
\1
=
\int \dd(\bar\eta,\eta)\,
\e^{-\bar\eta\eta}
\lvert \eta\rangle\langle \bar\eta\rvert,
\qquad
\bar\eta\eta\equiv \sum_i \bar\eta_i\eta_i,
\label{eq:coherent-identity}
\end{equation}
and the fermionic trace of an operator $O$ is
\begin{equation}
\Tr O
=
\int \dd(\bar\eta,\eta)\,
\e^{-\bar\eta\eta}
\langle -\bar\eta\rvert O\rvert \eta\rangle.
\label{eq:trace-formula}
\end{equation}

For $n$ replicas, introduce independent Grassmann variables
\begin{equation}
\eta^{(\alpha)}=(\eta_A^{(\alpha)},\eta_B^{(\alpha)}),
\qquad
\bar\eta^{(\alpha)}=(\bar\eta_A^{(\alpha)},\bar\eta_B^{(\alpha)}),
\qquad
\alpha=1,\dots,n,
\label{eq:replica-grassmann}
\end{equation}
and replicated coherent states
\begin{equation}
\lvert \eta\rangle
\equiv \bigotimes_{\alpha=1}^n \lvert \eta^{(\alpha)}\rangle,
\qquad
\langle \bar\eta\rvert
\equiv \bigotimes_{\alpha=1}^n \langle \bar\eta^{(\alpha)}\rvert.
\label{eq:replica-coherent}
\end{equation}

Applying \eqref{eq:trace-formula} to \eqref{eq:graded-swap-identity},
\begin{equation}
\widetilde Z_n[A]
=
\int \prod_{\alpha=1}^n \dd\br{\bar\eta^{(\alpha)},\eta^{(\alpha)}}
\exp\sq{-\sum_{\alpha=1}^n \bar\eta^{(\alpha)}\eta^{(\alpha)}}
\,
\langle -\bar\eta\rvert S_A^{\mathrm{gr}}
\br{\e^{-\beta H}}^{\otimes n}\rvert \eta\rangle.
\label{eq:Zn-before-path-int}
\end{equation}

The graded cyclic permutation modifies the coherent-state sewing on $A$ by the
Koszul sign. After the standard Grassmann bookkeeping, one obtains the mixed
boundary conditions
\begin{equation}
\psi_A^{(\alpha)}(\beta^-)=\psi_A^{(\alpha+1)}(0^+),
\qquad
\alpha=1,\dots,n-1,
\label{eq:bc-A-nonfinal}
\end{equation}
\begin{equation}
\psi_A^{(n)}(\beta^-)=-\psi_A^{(1)}(0^+),
\label{eq:bc-A-final}
\end{equation}
while on $B$ one keeps the ordinary anti-periodic closure replica by replica:
\begin{equation}
\psi_B^{(\alpha)}(\beta^-)
=
-\psi_B^{(\alpha)}(0^+),
\qquad
\alpha=1,\dots,n.
\label{eq:bc-B}
\end{equation}

Equivalently,
\begin{equation}
\widetilde Z_n[A]
=
\int \prod_{\alpha=1}^n \cD[\bar\psi^{(\alpha)},\psi^{(\alpha)}]
\exp\sq{-\sum_{\alpha=1}^n S[\bar\psi^{(\alpha)},\psi^{(\alpha)}]}
\label{eq:Zn-path-integral}
\end{equation}
with Euclidean action
\begin{equation}
S[\bar\psi,\psi]
=
\int_0^\beta \dd\tau\,
\cb{
\sum_i \bar\psi_i(\tau)\,\partial_\tau\psi_i(\tau)
+
H\br{\bar\psi(\tau),\psi(\tau)}
}
\label{eq:euclidean-action}
\end{equation}
and the graded mixed gluing \eqref{eq:bc-A-nonfinal}--\eqref{eq:bc-B}.

\section{Matrix form of the graded gluing and graded replica Fourier modes}

Let $P_A$ and $P_B$ be the one-particle projectors onto $A$ and $B$:
\begin{equation}
P_A+P_B=\1,
\qquad
P_A P_B=0.
\label{eq:AB-projectors}
\end{equation}
Introduce the replica-space vector
\begin{equation}
\Psi_i(\tau)
=
\begin{pmatrix}
\psi_i^{(1)}(\tau)\\
\psi_i^{(2)}(\tau)\\
\vdots\\
\psi_i^{(n)}(\tau)
\end{pmatrix}.
\label{eq:replica-vector}
\end{equation}

The graded shift matrix $K_n$ is the $n\times n$ matrix
\begin{equation}
K_n
\begin{pmatrix}
x_1\\x_2\\ \vdots\\ x_n
\end{pmatrix}
=
\begin{pmatrix}
x_2\\x_3\\ \vdots\\ -x_1
\end{pmatrix}.
\label{eq:graded-shift-action}
\end{equation}
Equivalently,
\begin{equation}
(K_n)_{\alpha\beta}
=
\delta_{\alpha,\beta-1}
-\delta_{\alpha,n}\delta_{\beta,1}.
\label{eq:graded-shift-matrix}
\end{equation}
Then the mixed boundary conditions become
\begin{equation}
\Psi(\beta^-)
=
\sq{
-P_B\otimes \1_n
+
P_A\otimes K_n
}\Psi(0^+).
\label{eq:matrix-bc-graded}
\end{equation}

Because $K_n^n=-\1_n$, its eigenvalues are the $n$ roots of $-1$:
\begin{equation}
\lambda_m=\e^{\ii q_m},
\qquad
q_m=\frac{(2m+1)\pi}{n},
\qquad
m=0,\dots,n-1.
\label{eq:graded-eigs}
\end{equation}
We therefore define the graded replica Fourier modes by
\begin{equation}
\psi_i^{(m)}(\tau)
=
\frac{1}{\sqrt n}
\sum_{\alpha=1}^n
\e^{-\ii q_m \alpha}\psi_i^{(\alpha)}(\tau).
\label{eq:graded-replica-FT}
\end{equation}
These diagonalize the graded shift:
\begin{equation}
K_n \psi^{(m)} = \e^{\ii q_m}\psi^{(m)}.
\label{eq:K-eigenvalue}
\end{equation}

Hence the sector boundary conditions are
\begin{equation}
\psi^{(m)}(\beta^-)
=
\br{
-P_B
+
\e^{\ii q_m}P_A
}\psi^{(m)}(0^+).
\label{eq:sector-bc-graded-direct}
\end{equation}
It is convenient to rewrite this as
\begin{equation}
\psi^{(m)}(\beta^-)
=
-\Omega_m \psi^{(m)}(0^+),
\qquad
\Omega_m\equiv P_B-\e^{\ii q_m}P_A.
\label{eq:Omega-def}
\end{equation}

Special cases:
\begin{itemize}[leftmargin=2em]
\item $A=\varnothing$: $\Omega_m=\1$, so one recovers the ordinary
anti-periodic boundary condition.
\item $A=\HH$: $\Omega_m=-\e^{\ii q_m}\1$, so the full system acquires the
correct fermionic half-shifted phases $\e^{\ii q_m}$.
\item $n=2$: $q_0=\pi/2$ and $q_1=3\pi/2$, so the two $A$-sector phases are
$+\ii$ and $-\ii$.
\end{itemize}

\section{Quadratic Hamiltonian and the convention for the free Green's function}

We now specialize to the Hubbard model
\begin{equation}
H
=
-\sum_{ij,\sigma} t_{ij}\, c_{i\sigma}^\dagger c_{j\sigma}
+
U\sum_i n_{i\up}n_{i\dn}
-
\mu \sum_{i,\sigma} n_{i\sigma}.
\label{eq:Hubbard-H}
\end{equation}
We split it into
\begin{equation}
H=H_0+H_{\mathrm{int}},
\qquad
H_{\mathrm{int}}=U\sum_i n_{i\up}n_{i\dn},
\label{eq:H-split}
\end{equation}
with one-body kernel
\begin{equation}
(H_0)_{ij,\sigma}=h_{ij}=-t_{ij}-\mu\delta_{ij}.
\label{eq:h-def}
\end{equation}
The quadratic Euclidean action is
\begin{equation}
S_0
=
\int_0^\beta \dd\tau \sum_{ij,\sigma}
\bar c_{i\sigma}(\tau)\sq{\delta_{ij}\partial_\tau + h_{ij}}c_{j\sigma}(\tau).
\label{eq:S0}
\end{equation}

\subsection*{Convention}
We fix the sign convention
\begin{equation}
\boxed{
\bar G_\sigma(X,X')
\equiv
\langle T_\tau\, c_\sigma(X)c_\sigma^\dagger(X')\rangle_0,
}
\label{eq:Gbar-def-single}
\end{equation}
with \emph{no} extra minus sign in front of $T_\tau$.

\begin{remark}
In many textbooks one instead defines
\[
G_\sigma^{\mathrm{std}}(X,X')\equiv
-\langle T_\tau\, c_\sigma(X)c_\sigma^\dagger(X')\rangle_0.
\]
The two conventions are related by
\[
\bar G_\sigma=-G_\sigma^{\mathrm{std}}.
\]
Therefore any formula imported from the standard convention must acquire an
overall minus sign when rewritten in terms of $\bar G$.
\end{remark}

\section{Free graded replica Green's function on the $n$-sheeted manifold}

\subsection{Definition in replica space}

For the replicated free theory with the graded boundary conditions
\eqref{eq:bc-A-nonfinal}--\eqref{eq:bc-B}, define
\begin{equation}
\bar G_{\sigma;\alpha\beta}^{(n)}(i,\tau;j,\tau')
\equiv
\left\langle
T_\tau\,
c_{i\sigma}^{(\alpha)}(\tau)\,
c_{j\sigma}^{(\beta)\dagger}(\tau')
\right\rangle_{0,n}.
\label{eq:Gbar-replica-def}
\end{equation}
Passing to the graded sector basis \eqref{eq:graded-replica-FT}, define
\begin{equation}
\bar G_\sigma^{(m)}(i,\tau;j,\tau')
\equiv
\left\langle
T_\tau\,
\psi_{i\sigma}^{(m)}(\tau)\,
\psi_{j\sigma}^{(m)\dagger}(\tau')
\right\rangle_{0,n}.
\label{eq:Gm-def}
\end{equation}
The inverse transform is
\begin{equation}
\bar G_{\sigma;\alpha\beta}^{(n)}(i,\tau;j,\tau')
=
\frac{1}{n}
\sum_{m=0}^{n-1}
\e^{\ii q_m(\alpha-\beta)}
\bar G_\sigma^{(m)}(i,\tau;j,\tau').
\label{eq:G-replica-from-m}
\end{equation}

\subsection{Equation of motion and jump condition}

Because the convention \eqref{eq:Gbar-replica-def} carries no extra prefactor
minus sign, the equation of motion is
\begin{equation}
\boxed{
\br{\partial_\tau+\bh}\,\bar G_\sigma^{(m)}(\tau,\tau')
=
\delta(\tau-\tau')\,\bI.
}
\label{eq:EOM-Gbar}
\end{equation}
Integrating across $\tau=\tau'$ yields the equal-time jump
\begin{equation}
\boxed{
\bar G_\sigma^{(m)}(\tau'^+,\tau')
-
\bar G_\sigma^{(m)}(\tau'^-,\tau')
=
\bI.
}
\label{eq:jump-Gbar}
\end{equation}
The sector boundary condition is
\begin{equation}
\bar G_\sigma^{(m)}(\beta^-,\tau')
=
-\Omega_m\,\bar G_\sigma^{(m)}(0^+,\tau').
\label{eq:G-boundary-sector}
\end{equation}

\subsection{Exact solution in one graded sector}

Fix $\tau'$. For $\tau\neq\tau'$, the homogeneous equation implies
\begin{equation}
\bar G_\sigma^{(m)}(\tau,\tau')
=
\e^{-\tau\bh}\,C_\pm^{(m)}(\tau'),
\qquad
\text{for } \tau\gtrless\tau',
\label{eq:piecewise-homogeneous}
\end{equation}
with two a priori different matrices $C_+^{(m)}(\tau')$ and
$C_-^{(m)}(\tau')$.

Using the boundary condition \eqref{eq:G-boundary-sector}, with
$\beta^- > \tau'$ and $0^+<\tau'$, gives
\begin{equation}
\e^{-\beta\bh} C_+^{(m)}(\tau')
=
-\Omega_m\, C_-^{(m)}(\tau').
\label{eq:boundary-Cpm}
\end{equation}
The jump condition \eqref{eq:jump-Gbar} gives
\begin{equation}
\e^{-\tau' \bh}
\br{C_+^{(m)}(\tau')-C_-^{(m)}(\tau')}
=
\bI.
\label{eq:jump-Cpm}
\end{equation}

Define
\begin{equation}
B_m\equiv \bI+\Omega_m^{-1}\e^{-\beta\bh}.
\label{eq:Bm-def}
\end{equation}
Whenever $B_m$ is invertible, define
\begin{equation}
A_m\equiv B_m^{-1}
=
\sq{\bI+\Omega_m^{-1}\e^{-\beta\bh}}^{-1}.
\label{eq:Am-def}
\end{equation}
Solving \eqref{eq:boundary-Cpm}--\eqref{eq:jump-Cpm} gives
\begin{equation}
C_+^{(m)}(\tau') = A_m \e^{\tau'\bh},
\qquad
C_-^{(m)}(\tau') = \br{A_m-\bI}\e^{\tau'\bh}.
\label{eq:Cpm-solution}
\end{equation}
Therefore
\begin{equation}
\boxed{
\bar G_\sigma^{(m)}(\tau,\tau')
=
\begin{cases}
\e^{-\tau\bh}\,A_m\,\e^{\tau'\bh},
& \tau>\tau',\\[6pt]
\e^{-\tau\bh}\,\br{A_m-\bI}\,\e^{\tau'\bh},
& \tau<\tau'.
\end{cases}
}
\label{eq:Gm-piecewise}
\end{equation}
Equivalently,
\begin{equation}
\boxed{
\bar G_\sigma^{(m)}(\tau,\tau')
=
\e^{-\tau\bh}
\sq{
\Theta(\tau-\tau')A_m
+
\Theta(\tau'-\tau)\br{A_m-\bI}
}
\e^{\tau'\bh}.
}
\label{eq:Gm-step-function}
\end{equation}

The formula \eqref{eq:Gm-piecewise} is the correct generic solution. It
depends on $\tau$ and $\tau'$ separately unless $\bh$ commutes with $A_m$.

\paragraph{Right boundary condition.}
Because $\Omega_m$ is unitary,
\begin{equation}
\Omega_m^\dagger=\Omega_m^{-1}.
\label{eq:Omega-unitary}
\end{equation}
One checks directly from \eqref{eq:Gm-piecewise} that
\begin{equation}
\bar G_\sigma^{(m)}(\tau,\beta^-)
=
-\bar G_\sigma^{(m)}(\tau,0^+)\Omega_m^\dagger.
\label{eq:right-boundary}
\end{equation}

\subsection{Explicit reconstruction for $n=2$}

For $n=2$ one has
\begin{equation}
q_0=\frac{\pi}{2},
\qquad
q_1=\frac{3\pi}{2},
\qquad
\e^{\ii q_0}=+\ii,
\qquad
\e^{\ii q_1}=-\ii.
\label{eq:n2-phases}
\end{equation}
Hence
\begin{equation}
\bar G_{\sigma;11}^{(2)}=\bar G_{\sigma;22}^{(2)}
=
\frac{\bar G_\sigma^{(0)}+\bar G_\sigma^{(1)}}{2},
\label{eq:n2-diag}
\end{equation}
\begin{equation}
\bar G_{\sigma;12}^{(2)}
=
-\frac{\ii}{2}\br{\bar G_\sigma^{(0)}-\bar G_\sigma^{(1)}},
\qquad
\bar G_{\sigma;21}^{(2)}
=
\frac{\ii}{2}\br{\bar G_\sigma^{(0)}-\bar G_\sigma^{(1)}}.
\label{eq:n2-offdiag}
\end{equation}
For a real one-body Hamiltonian $\bh$, the two sector Green's functions are
complex conjugates, so the replica-space Green's functions are real.

\subsection{Numerical construction}

The numerical construction must keep the two times separately.

\medskip
\textbf{Input:} one-body Hamiltonian matrix $\bh$, projectors $P_A$ and
$P_B=\bI-P_A$, inverse temperature $\beta$, replica number $n$.

\medskip
\textbf{Step 1:} compute
\begin{equation}
U(\tau)\equiv \e^{-\tau\bh},
\qquad
U(\beta)\equiv \e^{-\beta\bh}.
\label{eq:U-tau}
\end{equation}

\textbf{Step 2:} for each sector $m=0,\dots,n-1$, define
\begin{equation}
q_m=\frac{(2m+1)\pi}{n},
\qquad
\Omega_m=P_B-\e^{\ii q_m}P_A.
\label{eq:Omega-m-numerical}
\end{equation}

\textbf{Step 3:} form
\begin{equation}
B_m=\bI+\Omega_m^{-1}U(\beta),
\qquad
A_m=B_m^{-1},
\label{eq:Bm-numerical}
\end{equation}
whenever $B_m$ is invertible.

\textbf{Step 4:} evaluate
\begin{equation}
\bar G_\sigma^{(m)}(\tau,\tau')
=
\begin{cases}
U(\tau)A_mU(-\tau'), & \tau>\tau',\\[4pt]
U(\tau)\br{A_m-\bI}U(-\tau'), & \tau<\tau'.
\end{cases}
\label{eq:Gm-numerical-general}
\end{equation}

\textbf{Step 5:} transform back to replica indices using
\begin{equation}
\bar G_{\sigma;\alpha\beta}^{(n)}(i,\tau;j,\tau')
=
\frac{1}{n}\sum_{m=0}^{n-1}\e^{\ii q_m(\alpha-\beta)}
\bar G_\sigma^{(m)}(i,\tau;j,\tau').
\label{eq:G-back-to-replica}
\end{equation}

\textbf{Step 6:} for Monte Carlo sampling, use
\begin{equation}
\bar G_{\sigma;\alpha\beta}^{(n)}(i,\tau_p;j,\tau_q)
=
\br{\bar G_{\sigma;\alpha\beta}^{(n)}(\tau_p,\tau_q)}_{ij}.
\label{eq:Gpq-mc}
\end{equation}
In general there is no reduction to a lookup table depending only on the
wrapped time difference.

\subsection{Physical interpretation: why time-translation invariance is broken}

In the ordinary thermal theory, the Euclidean time circle is homogeneous: one
closes the contour by the same anti-periodic rule everywhere, so correlators
depend only on relative time up to anti-periodic wrapping.

In the entanglement replica theory, the seam at $\tau=0/\beta$ is a genuine
defect line:
\begin{itemize}[leftmargin=2em]
\item on $B$, the field returns to the same replica with the ordinary
anti-periodic sign;
\item on $A$, the field is redirected by the graded shift $K_n$.
\end{itemize}
Thus the propagator can detect whether the path from $\tau'$ to $\tau$ crosses
the entanglement seam. If $\bh$ mixes $A$ and $B$, this destroys imaginary-time
translation invariance.

\subsection{Existence of $A_m$, singular sectors, and zero modes}

The issue is whether
\begin{equation}
A_m=\sq{\bI+\Omega_m^{-1}\e^{-\beta\bh}}^{-1}
\label{eq:Am-existence}
\end{equation}
exists. This requires the invertibility of
\begin{equation}
B_m=\bI+\Omega_m^{-1}\e^{-\beta\bh}.
\label{eq:Bm-existence}
\end{equation}

\begin{proposition}[Singular sector criterion]
The following are equivalent:
\begin{enumerate}[leftmargin=2.2em]
\item $B_m$ is singular;
\item there exists a nonzero vector $v$ such that
\begin{equation}
\Omega_m^{-1}\e^{-\beta\bh}v=-v;
\label{eq:zero-mode-eigen}
\end{equation}
\item there exists a nontrivial homogeneous solution
\begin{equation}
\psi(\tau)=\e^{-\tau\bh}v
\label{eq:homogeneous-mode}
\end{equation}
of $(\partial_\tau+\bh)\psi=0$ satisfying
\begin{equation}
\psi(\beta^-)=-\Omega_m\psi(0^+).
\label{eq:homogeneous-twisted-bc}
\end{equation}
\end{enumerate}
\label{prop:zero-mode}
\end{proposition}

\begin{proof}
The equivalence of (1) and (2) is immediate from the definition of $B_m$.
Given (2), the function \eqref{eq:homogeneous-mode} obeys
\[
\psi(\beta^-)=\e^{-\beta\bh}v=-\Omega_m v=-\Omega_m\psi(0^+),
\]
so (3) follows. Conversely, any nontrivial homogeneous solution satisfying
\eqref{eq:homogeneous-twisted-bc} yields \eqref{eq:zero-mode-eigen} by setting
$v=\psi(0^+)$.
\end{proof}

Introduce the Euclidean monodromy operator
\begin{equation}
M_m\equiv -\Omega_m^{-1}\e^{-\beta\bh}.
\label{eq:Mm-def}
\end{equation}
Then $B_m$ is singular iff $M_m$ has eigenvalue $+1$.

\paragraph{Commuting case.}
If $[\bh,\Omega_m]=0$, choose simultaneous eigenvectors $u$ with
\begin{equation}
\bh u=\varepsilon u,
\qquad
\Omega_m u=\omega_m u.
\label{eq:commuting-eigenvectors}
\end{equation}
Then the singularity condition becomes
\begin{equation}
1+\e^{-\beta\varepsilon}\omega_m^{-1}=0.
\label{eq:singularity-commuting}
\end{equation}
On $B$, one has $\omega_m=1$, so \eqref{eq:singularity-commuting} has no
solution for real $\varepsilon$. On $A$, one has
$\omega_m=-\e^{\ii q_m}$, where $q_m$ is half-shifted as in
\eqref{eq:graded-eigs}; again there is no full-system boson-like periodic
sector. Thus the spurious full-system periodic zero mode that appears in the
naive bosonic-looking construction is absent in the corrected fermionic
graded-swap theory.

\section{Perturbation expansion and determinant structure}

We now return to the interacting Hubbard model \eqref{eq:Hubbard-H}. In the
replicated theory the interaction is diagonal in replica space:
\begin{equation}
H_{\mathrm{int}}^{(n)}
=
U\sum_{\alpha=1}^n\sum_i
n_{i\up}^{(\alpha)}n_{i\dn}^{(\alpha)}.
\label{eq:Hint-replica}
\end{equation}
Define the free graded-replica average $\langle\cdots\rangle_{0,n}$ using the
free replica Green's function \eqref{eq:Gbar-replica-def}. Then the
\emph{unnormalized twisted} object expands as
\begin{equation}
\widetilde Z_n[A]
=
\widetilde Z_{0,n}[A]\,
\left\langle
T_\tau \exp\sq{
-\int_0^\beta \dd\tau\,
U\sum_{\alpha=1}^n\sum_i
n_{i\up}^{(\alpha)}(\tau)n_{i\dn}^{(\alpha)}(\tau)
}
\right\rangle_{0,n}.
\label{eq:Zn-expand-start}
\end{equation}

Let
\begin{equation}
X_\ell\equiv (i_\ell,\tau_\ell,\alpha_\ell),
\qquad
\ell=1,\dots,k.
\label{eq:Xl-def}
\end{equation}
At order $k$ one finds
\begin{equation}
\frac{\widetilde Z_n^{(k)}[A]}{\widetilde Z_{0,n}[A]}
=
\frac{(-U)^k}{k!}
\sum_{i_1,\dots,i_k}
\sum_{\alpha_1,\dots,\alpha_k}
\int_0^\beta \dd\tau_1\cdots \int_0^\beta \dd\tau_k
\,
\det\sq{\bar G_{\up}^{(n)}(X_p,X_q)}_{p,q=1}^k
\,
\det\sq{\bar G_{\dn}^{(n)}(X_p,X_q)}_{p,q=1}^k.
\label{eq:Zn-k-det-full-domain}
\end{equation}
Equivalently, over ordered times
\begin{equation}
0<\tau_1<\tau_2<\cdots<\tau_k<\beta,
\label{eq:ordered-times}
\end{equation}
the factor $1/k!$ disappears:
\begin{equation}
\frac{\widetilde Z_n^{(k)}[A]}{\widetilde Z_{0,n}[A]}
=
(-U)^k
\sum_{i_1,\dots,i_k}
\sum_{\alpha_1,\dots,\alpha_k}
\int_{0<\tau_1<\cdots<\tau_k<\beta}
\prod_{\ell=1}^k \dd\tau_\ell\,
\det\sq{\bar G_{\up}^{(n)}(X_p,X_q)}_{p,q=1}^k
\,
\det\sq{\bar G_{\dn}^{(n)}(X_p,X_q)}_{p,q=1}^k.
\label{eq:Zn-k-det-ordered}
\end{equation}

Thus
\begin{equation}
\ln Z_n[A]
=
\ln \widetilde Z_n[A]-n\ln Z_1
=
\ln \widetilde Z_{0,n}[A]-n\ln Z_{0,1}
\,+\,\text{connected corrections}.
\label{eq:logZn-connected}
\end{equation}
The subtraction by $n\ln Z_1$ is not optional; it is part of the definition
of the normalized object $Z_n[A]=\Tr_A(\rho_A^n)$.

\section{Shifted expansion}

\subsection{Homotopic Hamiltonian and the correct sign of the shift}

For connected-determinant style reorganizations, introduce the homotopic
Hamiltonian
\begin{equation}
H(\xi)
=
-\sum_{ij,\sigma} t_{ij}\, c_{i\sigma}^\dagger c_{j\sigma}
+
\xi\br{
U\sum_i n_{i\up}n_{i\dn}
-
\sum_{i,\sigma}\alpha_\sigma n_{i\sigma}
}
-
\sum_{i,\sigma}(\mu_\sigma-\alpha_\sigma)n_{i\sigma}.
\label{eq:Hxi}
\end{equation}
Equivalently,
\begin{equation}
H(\xi)=H_0' + \xi\,H_{\mathrm{int}}',
\label{eq:Hxi-split}
\end{equation}
with
\begin{align}
H_0'
&=
-\sum_{ij,\sigma} t_{ij}\, c_{i\sigma}^\dagger c_{j\sigma}
-
\sum_{i,\sigma}(\mu_\sigma-\alpha_\sigma)n_{i\sigma},
\label{eq:H0prime}
\\
H_{\mathrm{int}}'
&=
U\sum_i n_{i\up}n_{i\dn}
-
\sum_i \alpha_\up n_{i\up}
-
\sum_i \alpha_\dn n_{i\dn}.
\label{eq:Hintprime}
\end{align}
This can be rewritten as
\begin{equation}
H_{\mathrm{int}}'
=
U\sum_i
\cb{
\br{n_{i\up}-\frac{\alpha_\dn}{U}}
\br{n_{i\dn}-\frac{\alpha_\up}{U}}
-
\frac{\alpha_\up\alpha_\dn}{U^2}
}.
\label{eq:Hint-shifted-factorized}
\end{equation}
Therefore the determinant shift is
\begin{equation}
\boxed{
\bar G_\sigma^{(n)}(X_p,X_q)
\ \longrightarrow\
\widetilde G_\sigma^{(n)}(X_p,X_q)
\equiv
\bar G_\sigma^{(n)}(X_p,X_q)
+
\delta_{pq}\,\frac{\alpha_{\bar\sigma}}{U}.
}
\label{eq:shifted-G}
\end{equation}
The shift is positive and occurs in the opposite spin channel.

\subsection{Shifted determinant structure}

With the shift \eqref{eq:shifted-G}, the twisted order-$k$ integrand takes the
form
\begin{equation}
\det\sq{\widetilde G_{\up}^{(n)}(X_p,X_q)}_{p,q=1}^k
\;
\det\sq{\widetilde G_{\dn}^{(n)}(X_p,X_q)}_{p,q=1}^k.
\label{eq:shifted-determinant}
\end{equation}
For $k\ge 2$, the artificial terms generated by the diagonal shifts are
disconnected and are removed in the connected expansion of
$\ln \widetilde Z_n[A]$.

\subsection{Why the first order is special}

At $k=1$, the diagonal shift would naively generate a term corresponding to
two crosses attached to the same interaction line. This is precisely the
constant subtraction in \eqref{eq:Hint-shifted-factorized}, so it must not be
double counted.

The first-order contribution to the \emph{twisted unnormalized} object is
\begin{equation}
\boxed{
\partial_\xi \ln \widetilde Z_n[A]\Big|_{\xi=0}
=
-U
\sum_{\alpha=1}^n
\sum_i
\int_0^\beta \dd\tau\,
\cb{
\br{
\bar G_{\up;\alpha\alpha}^{(n)}(i,\tau;i,\tau^+)
+\frac{\alpha_\dn}{U}
}
\br{
\bar G_{\dn;\alpha\alpha}^{(n)}(i,\tau;i,\tau^+)
+\frac{\alpha_\up}{U}
}
-
\frac{\alpha_\up\alpha_\dn}{U^2}
}.
}
\label{eq:lnZn1-twisted}
\end{equation}
This is the sign-correct and shift-correct first-order formula in the present
convention.

The equal-time symbol in \eqref{eq:lnZn1-twisted} is a lower-branch limit.
In \(\bar G(i,\tau;i,\tau^+)\), the second time argument is infinitesimally
later than the first one, so the piecewise formula must be evaluated with
\(\tau_1<\tau_2\):
\[
  \bar G^{(m)}(\tau,\tau^+)
  =
  \e^{-\tau\bh}\br{A_m-\bI}\e^{\tau\bh}.
\]
The upper branch would instead correspond to
\(\bar G^{(m)}(\tau^+,\tau)\).  In the C++ implementation this convention is
encoded by setting \(\tau_1=\tau_2=\tau\) as Chebyshev coordinates but passing
\(\texttt{use\_upper\_branch=false}\); the equality of the stored coordinates
represents the lower-side limit, not an average of the two branches.

For the normalized quantity $Z_n[A]=\widetilde Z_n[A]/Z_1^n$, one must also
subtract the ordinary one-replica thermal contribution:
\begin{equation}
\boxed{
\partial_\xi \ln Z_n[A]\Big|_{\xi=0}
=
\partial_\xi \ln \widetilde Z_n[A]\Big|_{\xi=0}
-
n\,\partial_\xi \ln Z_1\Big|_{\xi=0}.
}
\label{eq:lnZn1-normalized}
\end{equation}
Here
\begin{equation}
\partial_\xi \ln Z_1\Big|_{\xi=0}
=
-U
\sum_i
\int_0^\beta \dd\tau\,
\cb{
\br{
\bar G_{\up}^{\mathrm{th}}(i,\tau;i,\tau^+)
+\frac{\alpha_\dn}{U}
}
\br{
\bar G_{\dn}^{\mathrm{th}}(i,\tau;i,\tau^+)
+\frac{\alpha_\up}{U}
}
-
\frac{\alpha_\up\alpha_\dn}{U^2}
}.
\label{eq:lnZ1-first-order}
\end{equation}

For $n=2$, equation \eqref{eq:lnZn1-normalized} becomes
\begin{equation}
\partial_\xi \ln Z_2[A]\Big|_{\xi=0}
=
\partial_\xi \ln \widetilde Z_2[A]\Big|_{\xi=0}
-
2\,\partial_\xi \ln Z_1\Big|_{\xi=0}.
\label{eq:n2-normalized-first-order}
\end{equation}
This is exactly the object compared against the exact-diagonalization
finite-difference calculation.

\section{Relation to the entropy}

The normalized Renyi entropy is
\begin{equation}
S_A^{(n)}=\frac{1}{1-n}\ln Z_n[A]
=
\frac{1}{1-n}\br{\ln \widetilde Z_n[A]-n\ln Z_1}.
\label{eq:entropy-final}
\end{equation}
Thus the subtraction by $n\ln Z_1$ is part of the entropy itself whenever one
organizes the calculation through the twisted object $\widetilde Z_n[A]$.

\section{Summary of the final convention set}

For ease of reference, here is the corrected convention set.

\begin{enumerate}[leftmargin=2.2em]
\item The normalized object is
\[
Z_n[A]=\Tr_A(\rho_A^n),\qquad
\rho=\frac{\e^{-\beta H}}{Z_1},
\qquad
Z_n[A]=\frac{\widetilde Z_n[A]}{Z_1^n}.
\]

\item The operator inserted on $A$ is the \emph{graded} cyclic permutation
$S_A^{\mathrm{gr}}$, not the plain cyclic permutation.

\item The graded replica shift matrix has eigenvalues
\[
\lambda_m=\e^{\ii(2m+1)\pi/n}.
\]
Accordingly, the sector momenta are
\[
q_m=\frac{(2m+1)\pi}{n}.
\]

\item The sector boundary condition is
\[
\psi^{(m)}(\beta^-)= -\Omega_m \psi^{(m)}(0^+),
\qquad
\Omega_m=P_B-\e^{\ii q_m}P_A.
\]

\item The free Green's function is defined by
\[
\bar G=\langle T_\tau c c^\dagger\rangle_0
\]
with no extra minus sign, so
\[
(\partial_\tau+\bh)\bar G=\delta\,\bI
\]
and
\[
\bar G(\tau^+,\tau)-\bar G(\tau^-,\tau)=\bI.
\]

\item In one graded sector,
\[
\bar G^{(m)}(\tau,\tau')
=
\begin{cases}
\e^{-\tau\bh}A_m\e^{\tau'\bh}, & \tau>\tau',\\[4pt]
\e^{-\tau\bh}(A_m-\bI)\e^{\tau'\bh}, & \tau<\tau',
\end{cases}
\qquad
A_m=\sq{\bI+\Omega_m^{-1}\e^{-\beta\bh}}^{-1},
\]
whenever $A_m$ exists.

\item The replica-space Green's function is reconstructed with the
half-shifted phases,
\[
\bar G_{\alpha\beta}^{(n)}(\tau,\tau')
=
\frac1n\sum_{m=0}^{n-1}\e^{\ii q_m(\alpha-\beta)}\bar G^{(m)}(\tau,\tau').
\]
For $n=2$, this gives the $\pm \ii$ sector reconstruction
\eqref{eq:n2-diag}--\eqref{eq:n2-offdiag}.

\item The determinant expansion naturally computes
$\widetilde Z_n[A]$, not the normalized $Z_n[A]$.

\item In the shifted expansion, the first-order twisted contribution is
\eqref{eq:lnZn1-twisted}, while the normalized first-order correction is
\eqref{eq:lnZn1-normalized}.
\end{enumerate}

\bigskip
This completes the corrected derivation: starting from the fermionic
graded-swap representation of $\Tr_A(\rho_A^n)$, continuing through the
graded path integral and the corrected free replica Green's functions, and
ending with the determinant perturbation theory and shifted expansion in a
fully consistent fermionic sign convention.

\end{document}
